\section{固定证据下的条件绑定表示对照}
\label{sec:domain-evidence-coverage}

四来源评价中的支持页齐全而回答错误现象，提示检索后的阅读与对应关系仍需单独检验。本节另行建立八个来源组的固定评价，研究问题限定为：当小模型获得完全相同的结构化语义记录时，按主体、部件、属性和条件组织这些记录，能否改善答案与来源引用？本节不追加训练，也不将这一表示试验当作恢复轨迹 SFT 的迁移验证。

\subsection{来源、问题与等信息约束}

资料覆盖四组气象雷达与四组船用雷达，来自八家制造商，按预先声明的十三页 PDF 与十三段 HTML 范围整理为 502 条来源读法。多主体共用的记录不按型号复制，记录数不解释为独立事实数。另由未读取整理值或模型输出的 AI 作者从原始来源编写每组十二题，共 96 道中文问题，再由两位非作者逐题审查。题型包括简单属性 16 题、量值形式 24 题、限定属性 20 题、表格绑定 28 题及限定列举 8 题。该过程形成的是有来源约束的 AI 参考，尚不构成人工金标；历史同族暴露与基座预训练暴露未知的限制继续保留。

检索只接收题目，固定采用 BM25、同一词表与 top-4 页面；再将检索页面映射到其中全部记录，不按参考答案筛选。四道题未取得全部参考页，仍计入各方法的 96 题分母。运行前的保守必要性审查确认 50 题包含竞争读数，覆盖全部八组，其中 31 题需要双限定或多个量值对应。该门槛用于判断对照是否值得开展，不能视为算法已经有效。

平铺表示将每条记录的主体、部件、属性、量值形式、原单位、完整条件、技术注记与来源标识逐条写出；绑定表示仅对连续且前缀相同的记录，按“主体$\rightarrow$部件$\rightarrow$属性$\rightarrow$条件”共享字段，叶节点继承这些字段。二者不重排记录、不合并重复项、不补充新的推断。记检索所得的声明语义记录序列为 $E_{\mathrm{sem}}(q)$，两种解码结果满足
\[
 \operatorname{Dec}_{\mathrm{flat}}(q)
 = \operatorname{Dec}_{\mathrm{bound}}(q)
 = E_{\mathrm{sem}}(q).
\]
这里的等信息约束只针对声明语义记录。首版完整归档记录因过长而未运行；在后续测量前登记一次修订，将引文原文、偏移和文本哈希留在审计附件中。模型未访问这些附件；原引文可能包含未整理的信息，因此不宣称对整个原文无损。全部 8,434 次记录出现、192 次表示往返及 502 条记录的附件还原得到 CPU 校验。原文页路线保留原始文本，作为应用参照；它与结构化路线的额外信息、上下文长度及整理成本不同，不是纯粹的表示因果对照。

\subsection{固定执行与独立判分}

三路统一使用未经本项目适配的 Qwen2.5-3B-Instruct，原生 32,768 上下文、最多 768 个生成 token、batch size 为 1，采用 bfloat16、SDPA 和贪心解码。raw、flat、bound 最长输入分别为 4,325、30,952、27,449 token；输入证据未截断。三个作业各固定运行全部 96 题，不因分数重跑。原文路线有一条回答达到输出上限，原样保留计分。

联合正确要求正文完整正确、实际给定证据支持其断言、且引用合法并支持断言。正文正确、证据支持、引用支持、题级绑定错误和格式诊断分别记录；仅 JSON 合规或引用标识存在均不足以通过联合指标。两位非题目作者的 AI 审阅者分别审阅全部 288 条输出，不去重。平铺与绑定输入在审阅时解码为相同字段和顺序并隐藏方法标签；原文外观仍可能可辨，不宣称完全盲审或独立真人专家判定。

两位审阅者对全部 288 条正文正确、联合正确及题级绑定错误的判断一致；证据支持、引用、冗余与错误子类共在十九条输出中出现二十三项字段分歧，保留原判后逐项裁决。裁决者已参与构造方法映射，不称盲裁决。分项裁决没有改变上述三项主结果。额外合法但无关的引用按集合覆盖断言的口径记冗余而不新增逐引用精度门槛；这一边界解释发生在审阅开始后，另行留痕，不伪称所有边界均在输出前穷尽。

\begin{table}[htbp]
\centering\small
\caption{八来源覆盖评价的 AI 语义结果；每路分母为 96}
\label{tab:coverage-answer}
\CoverageAnswerTable
\end{table}

\begin{table}[htbp]
\centering\small
\caption{各来源组的联合正确数；每组每路分母为 12}
\label{tab:coverage-groups}
\resizebox{\linewidth}{!}{\CoverageGroupTable}
\end{table}

\subsection{收益分解及适用边界}

表\ref{tab:coverage-answer}中，条件绑定相对平铺记录的正文正确数由 73 提高至 80，配对为\CoverageAnswerWins{}胜\CoverageAnswerLosses{}负，净增\CoverageAnswerNet{}题；题级绑定错误由 17 降至 11。联合正确从 17 提高至 52，配对三十七胜两负、净增三十五题，且八个来源组均为正净增。因此达到运行前登记的净增至少五题、至少三个来源组正增、绑定错误不增加的继续投入门槛。这是本项目的实践标准，不是统计显著性标准。

这一联合收益不能全部解释为事实或条件理解的增强。事后分解显示，三十七个联合改善中，\CoverageJointAlreadyCorrectWins{}个原本两路正文都正确，改善主要表现为引用可用；另外\CoverageJointAnswerWins{}个伴随平铺正文错误转为绑定正文正确。平铺与绑定仍分别有 56 与 28 条正文正确而引用失败，显示记录标识与来源页标识的混淆仍是重要障碍。该分解是输出后的描述性分析，不是预先登记的机制消融。

在预先登记的五十道竞争题中，正文正确为平铺 37、绑定 44，绑定错误为 11、5；三十一道复杂题对应的正文正确为 24、28，绑定错误为 6、2。这与保留雷达资料中的主体、部件、模式和测量条件的研究动机相符，但八个相关来源组与单一模型的一次运行不足以确证普遍的领域推理提升。

原文路线正文正确为 60，低于绑定的 80；其联合正确却为 56，高于绑定的 52。故本轮没有建立条件绑定对原文 RAG 的全面优势。原文路线还省去本批结构化读法整理，不能只报告平铺基线较低的联合结果而忽略这一参照。

\subsection{统一引用解析的事后强对照}

为判断联合指标的差异有多少可以通过通用后处理消除，在上述输出和分数均已见的条件下，另行登记统一引用解析规则。这是事后敏感性分析，不替换原始主结果，也不称为新的模型预测。对于每道题，仅从模型实际获得的证据建立记录标识到来源页标识的映射。模型输出中的引用若精确等于一个给定记录标识，且该记录唯一对应一个给定来源页，则替换为该来源页标识；原本合法的来源页标识保持不变。未知、歧义、拼接或缺失的引用不猜测、不补齐、不删除，重复引用与次序保持不变。三路应用完全相同的规则。

程序只接受完整 JSON 对象或完整单一 JSON 代码块，只允许引用数组内的标识字段发生变化。它不读取参考答案，不改正文、断言、拒答或其他字段，亦不新增模型调用。全部 288 条原输出、既有评分与失败题均保留；哈希和逐字段校验确认输入、答案及变动范围。平铺与绑定分别有 49、16 条回答发生引用标识替换，原文路没有替换。两位非题目作者的 AI 审阅者分别复审这 65 条变动回答，依据实际给定的证据判断引用是否支持全部断言；标识合法不自动计为引用支持。审阅者此前已见过原始匿名输出，因此不宣称全新或完全盲审；未发生替换的引用评分及全部正文、证据支持和绑定错误评分继承原判。

\begin{table}[htbp]
\centering\small
\caption{统一精确引用解析前后的联合正确数；事后对照，每路分母为96}
\label{tab:coverage-citation}
\resizebox{\linewidth}{!}{\CoverageCitationTable}
\end{table}

采用统一解析后，绑定相对平铺的联合正确配对为\CoverageResolvedWins{}胜、\CoverageResolvedLosses{}负，净增\CoverageResolvedNet{}题，见表\ref{tab:coverage-citation}。这一结果应与原始净增三十五题并列解释，而不能把可由通用引用解析消除的差异全部归给领域语义理解。正文正确仍为平铺 73、绑定 80，绑定错误仍为 17、11；后处理没有产生新的事实回答收益。该对照评估的是表示与统一解析的组合，不识别条件组织、序列长度及注意力分配各自的独立作用，也不证明恢复轨迹训练在雷达领域迁移成功。

统一解析后六个来源组正净增、两个持平，没有来源组负净增；两位审阅者对全部六十五条变动引用的支持判定一致，不需要覆盖裁决。平铺与绑定仍有 18 与 12 条正文正确但未通过联合指标，因此剩余十三题优势也不能全部说成事实推理收益。绑定加统一解析在本批的联合正确为 68，高于原文加同一规则的 56；这一观察仍受原文额外信息、上下文长度及结构化整理成本差异限制，不等于普遍或同成本优于原文检索增强问答。

\begin{table}[htbp]
\centering\small
\caption{固定覆盖评价的 token 与 GPU 任务存续成本}
\label{tab:coverage-cost}
\CoverageCostTable
\end{table}

三路正式评价合计 0.147056 GPU 小时，含单卡子进程的导入、加载、生成和清理，不是 GPU 内核利用时间。条件绑定相对平铺输入 token 减少 9.62\%，输入与生成总 token 减少 9.58\%，但两条结构化路线总 token 都明显高于原文。三路各用一张卡的一次时长仅作实测记录，不作通用速度保证；来源整理、AI 审阅与历史训练成本未计入此表。

据此，将条件绑定保留为固定结构化证据下的领域辅助表示结果；它属于确定性证据组织策略，没有训练新模型，也没有新增图关系推理或单位计算。本轮不继续用这些题目改提示或搜索表示。后续优先完善外部来源核验与参考可信度，再在另行登记的新来源或模型设置上讨论复验；公开集恢复轨迹 SFT 仍是论文的主要训练方法证据。
